% ================================================================
%  CHAPTER 4: RESULTS AND DISCUSSION
% ================================================================
\chapter{Results and Discussion}

This chapter presents the complete experimental results of the five-phase
computational pipeline. Section~\ref{sec:data_dist} characterises the curated
EGFR bioactivity dataset. Section~\ref{sec:ml_perf} evaluates and compares all
seven classification models. Section~\ref{sec:vs_results} reports Nepali
medicinal plant virtual screening outcomes and lead prioritisation.
Section~\ref{sec:disc} provides an integrated discussion of all findings.

% ----------------------------------------------------------------
\section{Data Distribution Analysis}
\label{sec:data_dist}
% ----------------------------------------------------------------

\subsection{Dataset Composition and Class Balance}

The curated dataset comprised \textbf{10,523 unique EGFR-bioactive compounds}
retrieved from ChEMBL (CHEMBL203) following strict quality filtering. The binary
label distribution exhibited a moderate class imbalance: \textbf{7,792 active
compounds} (74.0\%, \texorpdfstring{pIC$_{50}$}{pIC50} $\geq 6.0$) and
\textbf{2,731 inactive compounds} (26.0\%,
\texorpdfstring{pIC$_{50}$}{pIC50} $< 6.0$). This imbalance is inherent to
bioactivity databases and reflects the practical reality that published compounds
are disproportionately those that demonstrated measurable activity in assay
screening campaigns.

Table~\ref{tab:dataset_overview} summarises the dataset composition across the
three scaffold-split partitions.

\begin{table}[H]
\centering
\caption{Dataset composition and class distribution across scaffold-split partitions.}
\label{tab:dataset_overview}
{\singlespacing
\begin{tabular}{lcccc}
\toprule
\textbf{Partition} & \textbf{Total} & \textbf{Active} & \textbf{Inactive}
                   & \textbf{Active (\%)} \\
\midrule
Training   & 8{,}418 & 6{,}292 & 2{,}126 & 74.7 \\
Validation & 1{,}052 &     743 &     309 & 70.6 \\
Test       & 1{,}053 &     757 &     296 & 71.9 \\
\midrule
\textbf{Total} & \textbf{10{,}523} & \textbf{7{,}792} & \textbf{2{,}731}
               & \textbf{74.0} \\
\bottomrule
\end{tabular}
}
\end{table}

The class balance remained broadly consistent across all three partitions
(74.7\% / 70.6\% / 71.9\%), with the held-out validation and test partitions being
marginally less active-enriched. This is expected under a strict Bemis-Murcko
scaffold split, in which the large shared scaffold families are concentrated in the
training set and the held-out partitions are dominated by novel singleton scaffolds.
The overall class balance is shown in Figure~\ref{fig:class_balance}.

\begin{figure}[htbp]
\centering
\includegraphics[width=0.5\textwidth]{eda_class_balance.pdf}
\caption{Class balance of the curated EGFR dataset: 74.0\% active and 26.0\%
         inactive compounds across 10,523 unique molecules.}
\label{fig:class_balance}
\end{figure}

\subsection{\texorpdfstring{pIC$_{50}$}{pIC50} Distribution}

The pIC$_{50}$ distribution across the full dataset exhibited a bimodal profile.
The active subpopulation was centred around $\text{pIC}_{50} \approx 7.5$--8.0
(corresponding to IC$_{50}$ values of 10--32 nM), reflecting the high-potency
compounds typically selected for further optimisation in medicinal chemistry
campaigns. The inactive subpopulation showed a peak near $\text{pIC}_{50} \approx
5.2$--5.5, representing weak binders at the micro-molar level. The binary threshold
at pIC$_{50} = 6.0$ (IC$_{50} = 1~\mu$M) was placed in a relatively low-density
valley between these two modes, providing a chemically meaningful separation rather
than an arbitrary cut-point. This is consistent with the standard medicinal chemistry
convention that defines a ``hit'' as a compound with measurable activity at or
below 1 $\mu$M. The bimodal structure and the placement of the decision threshold
are shown in Figure~\ref{fig:pic50}.

\begin{figure}[htbp]
\centering
\includegraphics[width=0.82\textwidth]{eda_pic50_distribution.pdf}
\caption{Distribution of EGFR bioactivity (pIC$_{50}$) for the curated dataset,
         coloured by activity label. The active and inactive subpopulations form
         two modes separated by the pIC$_{50} = 6.0$ decision threshold (dashed
         line, IC$_{50} = 1~\mu$M).}
\label{fig:pic50}
\end{figure}

\subsection{Physicochemical Property Analysis}

Analysis of the RDKit-computed physicochemical properties of the curated dataset
revealed that the majority of compounds occupied drug-like chemical space. Key
observations are summarised in Table~\ref{tab:physchem}.

\begin{table}[H]
\centering
\caption{Summary statistics of physicochemical properties of the curated dataset.}
\label{tab:physchem}
{\singlespacing
\begin{tabular}{lcccc}
\toprule
\textbf{Property} & \textbf{Min} & \textbf{Median} & \textbf{Mean} & \textbf{Max} \\
\midrule
Molecular Weight (Da) & 65.4  & 478.0  & 475.0  & 991.7 \\
LogP                  & $-$2.0 & 4.5   & 4.5    & 13.1  \\
H-Bond Donors         & 0      & 2      & 2.2    & 10    \\
H-Bond Acceptors      & 0      & 7      & 6.6    & 17    \\
TPSA (\AA$^{2}$)      & 0.0    & 96.3   & 95.2   & 273.0 \\
\bottomrule
\end{tabular}
}
\end{table}

The median molecular weight of 478 Da and median LogP of 4.5 sit close to the upper
limits of Lipinski's Rule of Five (Ro5), consistent with the physicochemical
characteristics of ATP-competitive kinase inhibitors, which frequently occupy the
upper boundary of drug-like chemical space. Approximately \textbf{47\%} of the
compounds were fully Lipinski-compliant (MW $\leq$ 500 Da, LogP $\leq$ 5,
HBD $\leq$ 5, HBA $\leq$ 10); the remainder most often exceed the LogP or molecular
weight thresholds, reflecting the lipophilic, higher-mass character typical of this
inhibitor class. The molecular weight and lipophilicity distributions are shown in
Figure~\ref{fig:properties}.

\begin{figure}[htbp]
\centering
\includegraphics[width=0.95\textwidth]{eda_property_distributions.pdf}
\caption{Distributions of molecular weight (left) and calculated LogP (right) for a
         random sample of the curated dataset. Dashed lines mark the Lipinski
         thresholds (MW = 500 Da, LogP = 5); most compounds cluster near these
         thresholds, characteristic of kinase inhibitors.}
\label{fig:properties}
\end{figure}

\subsection{Scaffold Diversity Analysis}

The Bemis-Murcko scaffold analysis identified \textbf{3,693 unique scaffolds}
across the full dataset. Of these, 2,418 (65.5\%) were singletons, comprising
only a single compound. The most frequently occurring scaffold, observed in
572 compounds, corresponds to the 4-anilinoquinazoline core, which is the
defining pharmacophore of first- and second-generation EGFR TKIs including
gefitinib, erlotinib, and lapatinib. Several further scaffolds each comprising
over one hundred compounds include the pyrrolopyrimidine core characteristic of
certain second-generation and third-generation inhibitors. This scaffold-level analysis
confirms that the dataset represents a structurally diverse yet pharmacophore-focused
chemical space. Figure~\ref{fig:scaffold} shows the most frequent scaffolds and the
overall split between singleton and multi-compound scaffolds, which motivates the
use of scaffold-based data splitting.

\begin{figure}[htbp]
\centering
\includegraphics[width=0.95\textwidth]{eda_scaffold_diversity.pdf}
\caption{Bemis-Murcko scaffold diversity of the curated dataset. Left: number of
         compounds sharing each of the most frequent scaffolds. Right: the split
         between singleton scaffolds (a single compound) and multi-compound
         scaffolds, illustrating the long-tailed scaffold distribution that
         necessitates scaffold-based splitting.}
\label{fig:scaffold}
\end{figure}

% ----------------------------------------------------------------
\section{Machine Learning Model Performance}
\label{sec:ml_perf}
% ----------------------------------------------------------------

\subsection{Individual Model Performance on the Scaffold Test Set}

All seven models were trained under identical scaffold split conditions and evaluated
on the held-out scaffold test set of 1,053 compounds. Results are presented in
Table~\ref{tab:results}. The classification threshold was set at $\hat{p} = 0.5$
for all threshold-dependent metrics (F1-Score, MCC, Accuracy).

\begin{table}[H]
\centering
\caption{Performance comparison of all models on the scaffold-held-out test set
         (threshold = 0.50 for F1, MCC, Accuracy).}
\label{tab:results}
{\singlespacing
\begin{tabular}{llccccc}
\toprule
\textbf{Model} & \textbf{Type} & \textbf{AUROC} & \textbf{AUPRC}
               & \textbf{F1} & \textbf{MCC} & \textbf{Accuracy} \\
\midrule
\rowcolor{tableblue}
Random Forest       & ML + ECFP4 & 0.8827 & 0.9496 & 0.8762 & 0.5072 & 0.8129 \\
MLP                 & ML + ECFP4 & 0.8547 & 0.9354 & 0.8629 & 0.5281 & 0.8053 \\
GCN                 & GNN        & 0.8454 & 0.9332 & 0.7877 & 0.4707 & 0.7312 \\
GraphSAGE           & GNN        & 0.8229 & 0.9179 & 0.8070 & 0.4706 & 0.7474 \\
GIN                 & GNN        & 0.8217 & 0.9163 & 0.8487 & 0.4865 & 0.7863 \\
GAT                 & GNN        & 0.8162 & 0.9176 & 0.7689 & 0.4473 & 0.7123 \\
AttentiveFP         & GNN        & 0.8087 & 0.9108 & 0.8254 & 0.3989 & 0.7521 \\
\bottomrule
\end{tabular}
}
\end{table}

Figure~\ref{fig:model_comparison} visualises these results across the four headline
metrics, making the relative ordering of the seven models directly comparable.

\begin{figure}[htbp]
\centering
\includegraphics[width=0.98\textwidth]{results_model_comparison.pdf}
\caption{Performance of all seven models on the scaffold-held-out test set across
         AUROC, AUPRC, F1-score, and MCC. The Random Forest and MLP fingerprint
         baselines lead on AUROC, while the graph neural networks form a competitive
         cluster led by GCN.}
\label{fig:model_comparison}
\end{figure}

\subsubsection{Random Forest Performance}

Random Forest achieved the highest test-set AUROC (0.8827) and AUPRC (0.9496)
among all models, together with the strongest F1-Score (0.8762); the MLP baseline
attained the highest MCC (0.5281), narrowly ahead of Random Forest (0.5072). The
strong performance of the two fingerprint-based baselines is consistent with
established observations in cheminformatics benchmarks: ensemble and shallow
neural methods built on ECFP4 fingerprints are highly competitive on bioactivity
datasets of this scale ($\sim$10,000 compounds), particularly when the training
and test chemical spaces share fingerprint-level structural similarity. The high
AUPRC value (0.9496) is especially significant given the 74:26 class imbalance,
indicating that Random Forest maintains high precision at high recall thresholds.

\subsubsection{Graph Neural Network Performance}

Among the five GNN architectures at default settings, \textbf{GCN achieved the
highest AUROC (0.8454)}, followed by GraphSAGE (0.8229), GIN (0.8217), GAT
(0.8162), and AttentiveFP (0.8087). Notably, the simpler convolution- and
sampling-based architectures outperformed the more elaborate attention-based
models on this metric---a common observation, since high-capacity,
molecule-specific architectures such as AttentiveFP are highly sensitive to
hyperparameter choices and tend to underperform simpler models until properly
tuned (Section~\ref{sec:hp_tuning}). The per-architecture characteristics are:
\begin{itemize}
    \item \textbf{GCN} provided the strongest and most stable default GNN, its
          degree-normalised convolution generalising well to novel scaffolds without
          extensive tuning.
    \item \textbf{GraphSAGE} performed comparably (AUROC 0.8229), with its
          concatenated self-and-neighbour update giving robust inductive
          generalisation on molecular graphs of this scale.
    \item \textbf{GIN} achieved the strongest F1-Score (0.8487) and MCC (0.4865)
          among the GNNs; its sum aggregation preserves neighbourhood multiset
          information that benefits threshold-dependent classification.
    \item \textbf{GAT} did not translate its learned attention weights into a higher
          AUROC than GCN at default settings, indicating that the additional
          parameters require tuning to pay off.
    \item \textbf{AttentiveFP}, despite its sophisticated two-level attention and
          native 12-dimensional edge-feature integration, recorded the lowest
          default AUROC (0.8087) and MCC (0.3989). This points to substantial
          untapped capacity that motivates the dedicated hyperparameter optimisation
          reported in Section~\ref{sec:hp_tuning}.
\end{itemize}

\subsubsection{GNN vs.\ Traditional ML Baselines}

The performance gap between the fingerprint baselines and all default GNN models
(AUROC difference: 0.037--0.074 relative to Random Forest) reflects a well-documented
phenomenon in molecular machine learning: GNNs require larger training datasets and
careful tuning to overcome the strong inductive bias provided by ECFP4 fingerprints.
Notably, the MLP baseline (AUROC 0.8547), which applies deep learning to the same
ECFP4 fingerprints, outperformed all five default GNN models, confirming that the
fixed fingerprint representation—rather than deep learning capacity alone—drives the
advantage of the baselines. The ECFP4 representation was refined through decades of
cheminformatics research specifically for molecular activity prediction, while GNN
architectures must learn to identify relevant substructures from first principles.
This gap is not closed by hyperparameter optimisation: tuning AttentiveFP
(Section~\ref{sec:hp_tuning}) improved its validation AUROC but left its test AUROC
essentially unchanged, so on the strict scaffold test set the fingerprint baselines
remained ahead of every graph neural network.

\subsection{Hyperparameter Tuning Results}
\label{sec:hp_tuning}

Bayesian hyperparameter optimisation was performed over 20 Optuna trials using the
Tree-structured Parzen Estimator sampler, with the objective of maximising
validation AUROC. The optimal hyperparameter configuration identified for
AttentiveFP is presented in Table~\ref{tab:optuna}.

\begin{table}[H]
\centering
\caption{Optimal AttentiveFP hyperparameters identified by Optuna Bayesian
         optimisation over 20 trials, with the resulting performance.}
\label{tab:optuna}
{\singlespacing
\begin{tabular}{lcc}
\toprule
\textbf{Hyperparameter} & \textbf{Default Value} & \textbf{Optimal Value} \\
\midrule
\texttt{hidden\_channels}  & 128   & 128                         \\
\texttt{num\_layers}       & 2     & 3                           \\
\texttt{num\_timesteps}    & 2     & 2                           \\
\texttt{dropout}           & 0.30  & 0.129                       \\
\texttt{batch\_size}       & 64    & 64                          \\
\texttt{learning\_rate}    & 0.001 & $2.61 \times 10^{-3}$       \\
\texttt{weight\_decay}     & 0.0   & $1.10 \times 10^{-5}$       \\
\midrule
\textbf{Best validation AUROC (20 trials)} & \multicolumn{2}{c}{\textbf{0.8632}} \\
\textbf{Tuned test AUROC} (cf.\ default 0.8087) & \multicolumn{2}{c}{0.7970} \\
\bottomrule
\end{tabular}
}
\end{table}

The optimisation raised the best \emph{validation} AUROC to \textbf{0.8632}, a
modest improvement that lifts AttentiveFP into the middle of the field on the
validation metric. On the held-out scaffold \emph{test} set, however, the tuned model
achieved an AUROC of only \textbf{0.7970}---essentially unchanged from the untuned
model (0.8087), still the lowest of the GNNs, and well below the Random Forest
(0.8827) and MLP (0.8547) fingerprint baselines. In other words, hyperparameter
tuning improves the validation-set fit but does not translate into better test-set
generalisation, and the gap to the fingerprint-based baselines is not closed on this
strict scaffold-split task---a realistic outcome consistent with the well-documented
difficulty of out-performing tuned fingerprint models on medium-sized bioactivity
datasets.

The selected hyperparameters are nonetheless informative:

\begin{itemize}
    \item \textbf{num\_layers = 3}: A third message-passing layer enables atoms to
          exchange information across up to three bonds, covering the spatial extent
          of most pharmacophoric motifs (the default used only two).
    \item \textbf{dropout = 0.129}: A lower dropout than the default (0.30) was
          preferred, indicating that the model benefits from retaining more of its
          representational capacity on this dataset.
    \item \textbf{learning\_rate = $2.61 \times 10^{-3}$}: A higher learning rate
          than the default ($10^{-3}$) accelerated convergence within the per-trial
          epoch budget.
    \item \textbf{weight\_decay = $1.10 \times 10^{-5}$}: A small non-zero L2 penalty
          (absent by default) provided mild regularisation against scaffold-specific
          overfitting.
    \item \textbf{hidden\_channels, num\_timesteps, and batch\_size} were retained at
          their default values (128, 2, and 64), indicating these were already
          near-optimal for the task.
\end{itemize}

% ----------------------------------------------------------------
\section{Nepali Medicinal Plant Virtual Screening Results}
\label{sec:vs_results}
% ----------------------------------------------------------------

\subsection{Screening Library Composition}

The Nepali medicinal plant screening library was assembled from the COCONUT 2.0
database by filtering for compounds annotated to seven traditional Nepali plant
species, then applying the same RDKit standardisation pipeline as the ChEMBL training
data. After salt removal, canonicalisation, the molecular-weight filter, and
deduplication by canonical SMILES, the final library comprised \textbf{535 unique
natural product compounds} distributed across the seven species as summarised in
Table~\ref{tab:library_comp}.

\begin{table}[H]
\centering
\caption{Composition of the Nepali medicinal plant screening library from COCONUT 2.0.}
\label{tab:library_comp}
{\singlespacing
\begin{tabular}{lcc}
\toprule
\textbf{Plant Species} & \textbf{Compound Class} & \textbf{Compounds} \\
\midrule
\textit{Nardostachys jatamansi} & Sesquiterpenes, coumarins    & 203 \\
\textit{Terminalia chebula}     & Hydrolysable tannins         & 195 \\
\textit{Tinospora cordifolia}   & Diterpenes, alkaloids        & 71 \\
\textit{Swertia chirayita}      & Xanthones, secoiridoids      & 23 \\
\textit{Berberis aristata}      & Isoquinoline alkaloids       & 19 \\
\textit{Rhododendron arboreum}  & Triterpenoids, flavonoids    & 15 \\
\textit{Ocimum sanctum}         & Phenylpropanoids, terpenoids & 9  \\
\midrule
\textbf{Total} & --- & \textbf{535} \\
\bottomrule
\end{tabular}
}
\end{table}

The library is dominated by \textit{Nardostachys jatamansi} (203 compounds) and
\textit{Terminalia chebula} (195), reflecting their extensive phytochemical
characterisation in COCONUT, whereas \textit{Ocimum sanctum} (9) and
\textit{Rhododendron arboreum} (15) are comparatively under-represented---an
illustration of the uneven database coverage of Nepali medicinal flora. The
distribution across the seven species is shown in Figure~\ref{fig:library_comp}.

\begin{figure}[htbp]
\centering
\includegraphics[width=0.9\textwidth]{screening_library_composition.pdf}
\caption{Composition of the Nepali medicinal plant screening library by species,
         totalling 535 unique natural product compounds retrieved from COCONUT 2.0.}
\label{fig:library_comp}
\end{figure}

\subsection{Multi-Model Screening Funnel and Hit Identification}

A preliminary screen scored the library with the tuned AttentiveFP model. However,
AttentiveFP was in fact the \emph{weakest} classifier on the scaffold test set
(Section~\ref{sec:ml_perf}). To avoid basing lead selection on the weakest model,
the prospective screen was re-run and deployed with the study's two
\emph{best-performing} models---the strongest graph neural network (GCN, test AUROC
0.8454) and the strongest overall classifier (Random Forest, test AUROC 0.8827)---applied
to the identical 535-compound library. Table~\ref{tab:vs_summary} compares the
progressive screening funnel obtained under all three models.

\begin{table}[H]
\centering
\caption{Nepali plant virtual screening funnel under the three models, showing the
         number of compounds retained at each activity and drug-likeness stage. The
         three models disagree sharply, revealing the model-sensitivity of the screen.}
\label{tab:vs_summary}
{\singlespacing
\begin{tabular}{lccc}
\toprule
\textbf{Filter Stage} & \textbf{AttentiveFP} & \textbf{GCN} & \textbf{Random Forest} \\
                      & \textbf{(0.81)}      & \textbf{(0.85)} & \textbf{(0.88)} \\
\midrule
Full plant compound library                     & 535 & 535 & 535 \\
Predicted active ($\hat{p} \geq 0.50$)          & 144 &  24 & 181 \\
Predicted active ($\hat{p} \geq 0.80$)          &  83 &   1 &   6 \\
Predicted active ($\hat{p} \geq 0.90$)          &  71 &   0 &   1 \\
Drug-like leads ($\hat{p} \geq 0.90$, Ro5, QED $\geq 0.40$) & 11 & 0 & 1 \\
\bottomrule
\end{tabular}
}
\end{table}

The three funnels differ dramatically, and this divergence is itself a central
result. The tuned AttentiveFP predicts 144 actives at $\hat{p} \geq 0.50$ and an
implausible 71 compounds at $\hat{p} \geq 0.90$; the best graph network (GCN)
predicts only 24 actives and \emph{not a single} compound at $\hat{p} \geq 0.90$;
and the best overall classifier (Random Forest) predicts 181 actives but only one
compound at $\hat{p} \geq 0.90$. The abundance of apparently high-confidence hits in
the AttentiveFP screen is therefore an artefact of poor probability calibration on
out-of-distribution natural products by the weakest model, and is \emph{not}
reproduced by either best-calibrated model. This model-sensitivity is the principal
caution of the screening campaign: virtual-screening hit lists on natural-product
chemical space depend strongly on the choice of classifier, and the most trustworthy
leads are those nominated by the strongest models---ideally by more than one. A
second, related artefact---over-confidence on large, non-drug-like natural
products---is also model-dependent in its severity: the weakest model (AttentiveFP)
scored large hydrolysable \textit{Terminalia chebula} tannins (e.g.,
pentagalloylglucose, MW $\approx 941$~Da) at $\hat{p} \approx 1.0$, and the best GNN
(GCN) assigned its highest scores to large \textit{Swertia} and \textit{Tinospora}
triterpenoids, all of which violate the Rule of Five. Only the best classifier
(Random Forest) placed drug-like compounds---the flavones apigenin and kaempferol---at
the very top of its ranking, relegating the tannins below them ($\hat{p} \approx
0.82$). In every case, non-drug-like large compounds still appear among the high
scorers, reinforcing the need to couple activity prediction with drug-likeness
filtering. The three funnels are compared in Figure~\ref{fig:funnel}.

\begin{figure}[htbp]
\centering
\includegraphics[width=0.9\textwidth]{screening_funnel.pdf}
\caption{Virtual screening funnel for the Nepali medicinal plant library under the
         three models (tuned AttentiveFP, best GNN = GCN, best classifier = Random
         Forest), showing the number of compounds retained at each prioritisation
         stage. The sharp disagreement between models---AttentiveFP retaining 71
         compounds at $\hat{p} \geq 0.90$ versus 0 for GCN and 1 for Random
         Forest---illustrates the model-sensitivity of the screen.}
\label{fig:funnel}
\end{figure}

\subsection{Consensus Leads from the Best-Performing Models}

Because no single classifier is authoritative, lead prioritisation is based on the
two best-performing models rather than the weak AttentiveFP screen. The top drug-like
(Lipinski-compliant) leads nominated by the Random Forest are listed in
Table~\ref{tab:leads_rf}, and those nominated by the GCN in Table~\ref{tab:leads_gcn}.

\begin{table}[H]
\small
\centering
\caption{Top drug-like (Lipinski-compliant) Nepali plant leads nominated by the
         best overall classifier (Random Forest), ranked by predicted activity.
         Listed compounds are those bearing common (trivial) names.}
\label{tab:leads_rf}
{\singlespacing
\begin{tabular}{clcccccc}
\toprule
\textbf{Rank} & \textbf{Compound} & \textbf{$\hat{p}$}
              & \textbf{MW (Da)} & \textbf{LogP} & \textbf{QED}
              & \textbf{Ro5} & \textbf{Plant} \\
\midrule
1 & Apigenin       & 0.9127 & 270.2 & 2.58 & 0.632 & $\checkmark$ & \textit{O.~sanctum}    \\
2 & Kaempferol     & 0.8340 & 286.2 & 2.28 & 0.546 & $\checkmark$ & \textit{N.~jatamansi}  \\
3 & Isogentisin    & 0.7282 & 258.2 & 2.37 & 0.655 & $\checkmark$ & \textit{S.~chirayita}  \\
4 & Luteolin       & 0.6789 & 286.2 & 2.28 & 0.511 & $\checkmark$ & \textit{T.~chebula}    \\
5 & Tinosporaside  & 0.6607 & 492.5 & 0.63 & 0.440 & $\checkmark$ & \textit{T.~cordifolia} \\
6 & Columbin       & 0.6180 & 358.4 & 2.53 & 0.613 & $\checkmark$ & \textit{T.~cordifolia} \\
7 & Palmarin       & 0.6042 & 374.4 & 1.74 & 0.591 & $\checkmark$ & \textit{T.~cordifolia} \\
8 & Nardostachysin & 0.5836 & 430.5 & 2.95 & 0.526 & $\checkmark$ & \textit{N.~jatamansi}  \\
\bottomrule
\end{tabular}
}
\end{table}

\begin{table}[H]
\small
\centering
\caption{Top drug-like (Lipinski-compliant) Nepali plant leads nominated by the best
         graph neural network (GCN), ranked by predicted activity. The GCN is
         markedly more conservative: no compliant compound exceeds $\hat{p} = 0.66$,
         and none reaches the $\hat{p} \geq 0.80$ threshold.}
\label{tab:leads_gcn}
{\singlespacing
\begin{tabular}{clcccccc}
\toprule
\textbf{Rank} & \textbf{Compound} & \textbf{$\hat{p}$}
              & \textbf{MW (Da)} & \textbf{LogP} & \textbf{QED}
              & \textbf{Ro5} & \textbf{Plant} \\
\midrule
1 & Rulepidadiene~A   & 0.6531 & 202.3 & 4.19 & 0.552 & $\checkmark$ & \textit{N.~jatamansi}  \\
2 & Nardostachysin    & 0.5915 & 430.5 & 2.95 & 0.526 & $\checkmark$ & \textit{N.~jatamansi}  \\
3 & Palmarin          & 0.5532 & 374.4 & 1.74 & 0.591 & $\checkmark$ & \textit{T.~cordifolia} \\
4 & Luteolin          & 0.5138 & 286.2 & 2.28 & 0.511 & $\checkmark$ & \textit{T.~chebula}    \\
5 & Jatamanvaltrate~H & 0.4995 & 442.5 & 1.70 & 0.405 & $\checkmark$ & \textit{N.~jatamansi}  \\
\bottomrule
\end{tabular}
}
\end{table}

Two observations stand out. First, the best models converge on \emph{flavones}---%
apigenin, luteolin and kaempferol---as the most drug-like predicted-active scaffolds
(Table~\ref{tab:flavones}). Flavones are a chemically coherent and literature-supported
class of natural ATP-competitive kinase binders, lending biological plausibility that
the AttentiveFP depside/coumarin hits lacked. Second, although the best classifier
(Random Forest) is comparatively well calibrated---its two highest-scoring compounds
are the drug-like flavones apigenin and kaempferol---the best GNN (GCN) still assigns
its highest raw probabilities to large, non-drug-like \textit{Swertia} and
\textit{Tinospora} triterpenoids, so drug-likeness filtering remains essential.

\begin{table}[H]
\centering
\caption{Predicted activity of the flavone scaffolds under the two best models. Only
         luteolin is predicted active ($\hat{p} \geq 0.50$) by \emph{both} models.}
\label{tab:flavones}
{\singlespacing
\begin{tabular}{lcccccl}
\toprule
\textbf{Flavone} & \textbf{RF $\hat{p}$} & \textbf{GCN $\hat{p}$}
                 & \textbf{MW (Da)} & \textbf{QED} & \textbf{Ro5} & \textbf{Plant} \\
\midrule
Apigenin   & 0.913 & 0.410 & 270.2 & 0.632 & $\checkmark$ & \textit{O.~sanctum}   \\
Luteolin   & 0.679 & 0.514 & 286.2 & 0.511 & $\checkmark$ & \textit{T.~chebula}   \\
Kaempferol & 0.834 & 0.320 & 286.2 & 0.546 & $\checkmark$ & \textit{N.~jatamansi} \\
\bottomrule
\end{tabular}
}
\end{table}

Figure~\ref{fig:leads_scatter} positions the Random Forest drug-like leads in the
predicted-activity versus drug-likeness plane.

\begin{figure}[htbp]
\centering
\includegraphics[width=0.92\textwidth]{screening_leads_scatter.pdf}
\caption{Random Forest drug-like Nepali plant lead candidates in the
         predicted-activity versus drug-likeness (QED) plane. All shown candidates are
         Lipinski-compliant and marker size is proportional to molecular weight.
         Dashed lines mark the $H_2$ thresholds ($\hat{p}$~$\geq$~0.80,
         QED~$\geq$~0.40); apigenin and kaempferol clear both.}
\label{fig:leads_scatter}
\end{figure}

\subsubsection{Highest-Confidence Drug-Like Lead: Apigenin (\textit{Ocimum sanctum})}

The single highest-confidence Lipinski-compliant lead is \textbf{apigenin} from
\textit{Ocimum sanctum} (holy basil, Tulsi), scored by the best overall classifier
(Random Forest) at $\hat{p} = 0.913$:
\begin{itemize}
    \item \textbf{Predicted activity}: $\hat{p} = 0.913$ (Random Forest; the single
          library compound exceeding $\hat{p} \geq 0.90$)
    \item \textbf{Molecular weight}: 270.2 Da (well within the Ro5 limit of 500 Da)
    \item \textbf{LogP}: 2.58 (favourable lipophilicity)
    \item \textbf{QED}: 0.632 (clearly drug-like)
    \item \textbf{Lipinski Ro5}: Fully compliant --- no violations
    \item \textbf{Source}: \textit{Ocimum sanctum} (Tulsi) --- a sacred and widely
          used plant in Nepali and South Asian traditional medicine
\end{itemize}
Apigenin is a well-characterised dietary flavone with documented anti-proliferative
and kinase-modulating activity, providing biological plausibility for further
investigation. It is, however, a \emph{model-dependent} lead: the best graph network
scores the same compound at only $\hat{p} = 0.410$, illustrating the model-sensitivity
established above.

\subsubsection{Top Consensus Lead: Luteolin}

The most robust nomination is \textbf{luteolin}, the only small, drug-like flavone
predicted active ($\hat{p} \geq 0.50$) by \emph{both} best models---Random Forest
($\hat{p} = 0.679$) and GCN ($\hat{p} = 0.514$)---with MW $= 286.2$~Da and QED
$= 0.511$. Because it is endorsed independently by the best fingerprint model and the
best graph model, luteolin represents the highest-priority \emph{consensus} candidate
for experimental EGFR inhibition assays, notwithstanding its sub-0.80 probabilities.

\subsection{Structural Diversity of Plant-Derived Leads}

Across the two best models the drug-like leads span several distinct chemotypes:
flavones (apigenin, luteolin, kaempferol), xanthones (isogentisin), clerodane
furanoditerpenoids (columbin, palmarin, tinosporaside), iridoid esters
(jatamanvaltrate~H, nardostachysin), and sesquiterpenoids (rulepidadiene~A). This
diversity suggests that several distinct binding modes to the EGFR ATP cleft may be
accessible from Nepali plant compounds; the recurrence of the planar, hydroxylated
flavone scaffold across the independent predictions of both best models is the
strongest single structural signal, and a richer starting point than a structurally
homogeneous synthetic inhibitor library.

% ----------------------------------------------------------------
\section{Discussion}
\label{sec:disc}
% ----------------------------------------------------------------

\subsection{Evaluation of the Research Hypotheses}

The two research hypotheses formulated in Chapter~1 are evaluated
against the results reported above.

\textbf{$H_1$ (Performance Hypothesis) --- not supported.} $H_1$ predicted that the
graph neural networks, with rich atom and bond featurisation and Bayesian
hyperparameter optimisation, would achieve superior AUROC to the Random Forest and
MLP fingerprint baselines under the Bemis-Murcko scaffold split. The results do not
support this hypothesis: the fingerprint baselines achieved the highest scaffold
test-set AUROC (Random Forest 0.8827, MLP 0.8547), and no GNN matched them---GCN was
the strongest GNN at 0.8454, and the tuned AttentiveFP reached a validation AUROC of
only 0.8632 with a test AUROC of just 0.7970. $H_1$ is therefore
\textbf{rejected} for this dataset: on a medium-sized, scaffold-split EGFR
bioactivity set, well-tuned fingerprint models remained more accurate than the graph
neural networks. This is itself an informative, honestly reported result, consistent
with prior observations that ECFP4-based models are strong baselines on datasets of
this scale.

\textbf{$H_2$ (Virtual Screening Hypothesis) --- partially supported.} $H_2$ predicted
that the screening pipeline would identify at least one fully Lipinski-compliant
natural-product lead with predicted activity $\hat{p} \geq 0.80$ and QED $\geq 0.40$.
The verdict is \emph{model-dependent}. The best overall classifier (Random Forest)
identifies two such compounds---apigenin ($\hat{p} = 0.913$, QED $= 0.632$) and
kaempferol ($\hat{p} = 0.834$, QED $= 0.546$)---so $H_2$ is \textbf{supported by the
best classifier}. However, the best graph network (GCN) identifies \emph{no}
Lipinski-compliant compound above $\hat{p} = 0.80$, so $H_2$ is \textbf{not supported
by the best GNN} at that confidence. Only when the threshold is relaxed to
$\hat{p} \geq 0.50$ do the two best models produce a consensus drug-like lead
(luteolin). $H_2$ is therefore reported as \textbf{partially supported}: the pipeline
does surface actionable, drug-like natural-product leads (flavones), but their
predicted confidence and their identity depend strongly on the model, and the
inflated confidence of the earlier AttentiveFP screen is not reproduced by the
best-calibrated models.

\subsection{Performance of Graph Neural Networks Relative to Baselines}

The results demonstrate that the relationship between GNN and traditional ML
performance is nuanced. Across both the default and tuned configurations, the
fingerprint-based baselines outperformed all GNN models on the scaffold test set
(Random Forest AUROC 0.8827 and MLP 0.8547, versus 0.8087--0.8454 for the GNNs).
This finding is consistent with the molecular machine learning literature and
reflects two factors: (i) ECFP4 fingerprints encode decades of human cheminformatics
knowledge through their circular neighbourhood construction, providing a strong
structural inductive bias; and (ii) medium-sized bioactivity datasets ($\sim$10,000
compounds) often favour sample-efficient ensemble methods over parameter-intensive
deep learning approaches.

Bayesian hyperparameter optimisation raised AttentiveFP's validation AUROC to
0.8632 (mid-field on the validation metric), but its tuned scaffold test AUROC
(0.7970) remained the lowest of the GNNs and well below the fingerprint baselines. On this dataset, therefore, the GNNs offered no accuracy advantage: the
Random Forest was the strongest classifier and, in the prospective screen, also the
better-calibrated scoring model. The value of the graph approach here is not superior
accuracy but a genuinely independent representation---operating directly on molecular
graphs with the 29-dimensional atomic node features and 12-dimensional edge features
that are invisible to ECFP4 fingerprints---which lets the best GNN (GCN) serve as a
complementary second opinion to the Random Forest and thereby supports the cross-model
consensus prioritisation used in the screen that follows.

\subsection{Significance of Scaffold-Based Evaluation}

A critical methodological contribution of this thesis is the strict application of
Bemis-Murcko scaffold splitting. This approach fundamentally differs from random
splitting, which can lead to optimistic performance estimates when structurally
similar molecules appear in both training and test sets. Under scaffold splitting,
the test set is composed entirely of compound-scaffold pairs whose core frameworks
were never seen during training, providing a stringent evaluation of the model's
ability to generalise to genuinely novel chemical matter.

\subsection{Ethnopharmacological Rationale for Natural Product Screening}

The hit rate observed in the Nepali plant screening was itself model-dependent,
ranging from 4.5\% (GCN) through 33.8\% (Random Forest) at $\hat{p} \geq 0.50$. Even
so, both best models enrich for the same drug-like chemotype---the \emph{flavones}
apigenin, luteolin and kaempferol---which reflects the ethnopharmacological
enrichment inherent in traditional medicinal plant selection: plants used for
centuries in the treatment of inflammatory and proliferative conditions are
biologically pre-filtered for bioactive chemical space. The prominence of the
hydroxylated flavone scaffold among the top compliant hits is consistent with the
established literature on plant-derived kinase modulators, as planar polyphenolic
flavones are well documented ATP-competitive kinase binders that share structural
features with the core pharmacophores of synthetic EGFR inhibitors
\parencite{wu2018moleculenet}. At the same time, both the weakest model (AttentiveFP)
and the best GNN (GCN) still rank large, non-drug-like polyphenols and triterpenoids
among their highest-scoring compounds---a reminder that predictions on chemical space
distant from the training distribution must be interpreted with caution and coupled
with drug-likeness filtering.

\subsection{Lead Compound Prioritisation Strategy}

The prioritisation strategy applied three criteria---high predicted activity,
Lipinski Rule of Five compliance, and consensus between the two best models---reflecting
target engagement potential, oral bioavailability, and robustness to the
model-sensitivity documented above. Under this strategy apigenin emerges as the
highest-confidence single-model lead and luteolin as the top cross-model consensus
lead. Critically, these compounds are derived from plants (\textit{Ocimum sanctum},
\textit{Terminalia chebula}) that are inexpensive, culturally familiar, and locally
accessible throughout Nepal, offering a potential pathway to affordable
EGFR-targeted therapy candidates without the high synthesis costs of novel synthetic
compounds.

It is important to note that the prioritisation in this study represents an
\textit{in silico} hypothesis, not experimental confirmation. The identified leads
require: (i) \textit{in vitro} EGFR kinase inhibition assay; (ii) cellular
antiproliferative testing in EGFR-mutant NSCLC cell lines (e.g., PC9, HCC827);
and (iii) selectivity profiling against a kinase panel.

\subsection{Limitations of the Current Results}

Several limitations of the results presented in this chapter should be acknowledged:

\begin{itemize}
    \item \textbf{Model-sensitivity of the screen}: The single most important
          limitation is that the virtual-screening hit list depends strongly on the
          choice of classifier (144 predicted actives for AttentiveFP versus 24 for
          GCN and 181 for Random Forest at $\hat{p} \geq 0.50$; 71 versus 0 versus 1
          at $\hat{p} \geq 0.90$). No single model can be treated as authoritative,
          and the predicted probabilities are not calibrated on out-of-distribution
          natural products. Consensus across the best models mitigates but does not
          eliminate this uncertainty; only experimental assay can resolve it.
    \item \textbf{Validation AUROC vs.\ Test AUROC}: The tuned AttentiveFP's
          validation AUROC (0.8632) exceeded its scaffold test AUROC ($\approx$0.80),
          a gap that reflects optimisation against the validation partition. Future
          work should employ a nested cross-validation scheme for unbiased test-set
          estimation.
    \item \textbf{Assay heterogeneity}: Training labels aggregate IC$_{50}$ values
          from hundreds of laboratories with different protocols, imposing a label
          noise ceiling on scaffold-split test-set performance.
    \item \textbf{ADMET-proxy limitations}: The physicochemical properties used
          do not capture metabolic stability, P-glycoprotein efflux, plasma protein
          binding, or \textit{in vivo} pharmacokinetics.
    \item \textbf{COCONUT 2.0 annotation completeness}: Organism taxonomy annotations
          in COCONUT 2.0 may be incomplete for some Nepali plant species, potentially
          missing additional bioactive compounds from these plants.
    \item \textbf{No 3D structural validation}: The absence of molecular docking
          against EGFR crystal structures means no binding pose or free energy
          estimate is available to corroborate model predictions for the top leads.
\end{itemize}